\subsection{La foto de ejemplo: el casco de referencia}\label{sec:casco}

Para que la matemática no quede en el aire, cada paso de la metodología se resuelve a continuación, dentro del mismo paso, con una fotografía real y con
los números que produce el programa: el casco que el profesor pidió usar como objeto de prueba (figura~\ref{fig:casco-entrada}).
La foto se subió con el botón \textbf{Subir foto} de la página (procesamiento bajo) y los valores de esta sección
salen de ejecutar cada paso por separado (\texttt{experimentos/casco\_paso\_a\_paso.py}); el modelo que aparece en la
pestaña del escáner (figura~\ref{fig:visor-bajo}) coincide con ellos, con $N=65\,314$ vértices y $M=130\,624$ caras.

\begin{figure}[htbp]
\centering
\includegraphics[height=8.2cm]{casco_entrada.jpg}
\hfill
\includegraphics[height=8.2cm]{casco_recorte.jpg}
\caption{Foto de referencia reducida a $960\times1280$ píxeles, con los ejes $u$ (columna) y $v$ (fila), el recuadro de
encuadre, el centro de la imagen $(c_x,c_y)$ y los tres píxeles A, B y C que se siguen en las cuentas (izquierda); objeto
recortado por la máscara (derecha).}
\label{fig:casco-entrada}
\end{figure}
